\documentclass{beamer}
\usepackage[T1]{fontenc}
\usepackage[english]{babel}
\usepackage{lmodern}
\usepackage{amsmath,amssymb,mathtools}
\usepackage{algorithm}
\usepackage{algorithmic}
\usepackage{graphicx}
\usepackage{subcaption}
\usepackage{caption}
\usepackage{tikz}
\usetikzlibrary{shapes.geometric,positioning,calc,fit}
\usetikzlibrary{overlay-beamer-styles}
\usepackage{booktabs,multicol,stackengine}
\usepackage[table,dvipsnames]{xcolor}
\usepackage{cancel}
\usepackage{moresize}
\usepackage{hyperref}

\usepackage{listings}
\definecolor{nus-orange}{RGB}{239,124,0}
\definecolor{nus-blue}{RGB}{0,61,124}
\definecolor{halfgray}{gray}{0.55}

\lstset{
  basicstyle=\ttfamily\footnotesize,
  keywordstyle=\bfseries\color{nus-blue},
  stringstyle=\color{ForestGreen},
  numbers=left,
  numberstyle=\scriptsize\color{halfgray},
  rulesepcolor=\color{nus-orange},
  frame=shadowbox,
  columns=flexible,
  showstringspaces=false,
  keepspaces=true,
  tabsize=2,
  breaklines=true
}

\renewcommand{\figurename}{Figure}
\renewcommand{\algorithmname}{Algorithm}

\usepackage{SHU}
% VITAL STUFF
\renewcommand\epsilon\varepsilon
\renewcommand\phi\varphi
\renewcommand\leq\leqslant
\renewcommand\geq\geqslant
\newcommand\defeq{\stackrel{\mathrm{def}}{=}}

% For revision
\newcommand\rev[1]{{#1}}

% COMMENTS
\newcommand\mikael[1]{\textcolor{orange}{[\textbf{Mikaël:} #1]}}
\newcommand\pierre[1]{\textcolor{magenta}{[\textbf{Pierre:} #1]}}
\newcommand\pratik[1]{\textcolor{blue}{[\textbf{Pratik:} #1]}}

% SMOL SYMBOLS
\newcommand\calF{\mathcal{F}}
\newcommand\RR{\mathbb{R}}
\newcommand\QQ{\mathbb{Q}}
\newcommand\bD{\mathbf{D}}

% REDUCTIONS
\newcommand{\tr}{\leq_{\mathsf{P}}}
\newcommand{\equivt}{\equiv_{\mathsf{P}}}

% PROBLEMS
\newcommand\escore{\mathsf{EScore}}
\newcommand\score{\mathsf{Score}}
\newcommand\eshapley{\escore_{\cshapley}}
\newcommand\shapley{\score_{\cshapley}}
\newcommand\ebanzhaf{\escore_{\cbanzhaf}}
\newcommand\banzhaf{\score_{\cbanzhaf}}
\newcommand\ev{\mathsf{EV}}
\newcommand\evs{\mathsf{EV}_\star}
\newcommand\ennv{\mathsf{ENV}} % \env was already defined
\newcommand\envss{\mathsf{ENV}_{\star,\star}}
\newcommand\mc{\mathsf{MC}}
\newcommand\env{\mathsf{ENV}}

% NOTATIONS
\newcommand\vars{\mathsf{Vars}}
\newcommand\out{\mathsf{out}} % for output gate
\newcommand\cshapley{c_{\mathrm{Shapley}}}
\newcommand\cbanzhaf{c_{\mathrm{Banzhaf}}}
\newcommand\ar{\mathsf{ar}}
\newcommand\const{\mathsf{Const}}
\newcommand{\pqe}{\mathrm{PQE}}
\newcommand\prov{\mathrm{Prov}}

% For algorithm2e
\newcommand{\hrulealg}[0]{\vspace{1mm} \hrule \vspace{1mm}}

% For alignment of numbers
\newcommand{\0}{\phantom{0}}

% Example
\newcommand{\phiex}{\phi_{\mathrm{ex}}}


% Relational algebra operators
\newcommand\proj[1]{\pi_{[\text{#1}]}}
\newcommand\sel[1]{\sigma_{[\text{#1}]}}
\newcommand\ren[2]{\rho(\text{#1},\text{#2})}
\newcommand\njoin[1]{\bowtie_{[\text{#1}]}}
\newcommand\AND{\wedge}
\newcommand\OR{\vee}

\title{IT5008: Tutorial 5 --- Relational Algebra}
\author{\href{https://pratik2358.github.io/}{Pratik Karmakar}}
\institute{School of Computing,\\ National University of Singapore}
\date{AY26/27 S1}

\begin{document}

\begin{frame}
  \titlepage
  \IfFileExists{nus-logo.png}{
    \begin{figure}[htpb]\centering
      \includegraphics[keepaspectratio, scale=0.18]{nus-logo.png}
    \end{figure}
  }{}
\end{frame}

\section{Scenario}

\begin{frame}{The Book Exchange Schema (Recap)}
\small
Same NUN Book Exchange scenario and schema as Tutorials 3 and 4.
\begin{itemize}
  \item \textbf{department}(\underline{department}, faculty)
  \item \textbf{student}(name, \underline{email}, year, department \emph{(FK)}, graduate)
  \item \textbf{book}(title, format, pages, language, authors, publisher, year, edition, ISBN10 \emph{(unique)}, \underline{ISBN13})
  \item \textbf{copy}(\underline{owner}, \underline{book}, \underline{copy}) --- FKs to \texttt{student} and \texttt{book}
  \item \textbf{loan}(borrower, owner, book, copy, borrowed, returned)
\end{itemize}
\medskip
\begin{block}{New this tutorial}
Writing queries as \textbf{relational algebra} expressions, and translating them to/from SQL --- including \textbf{universal quantification} without \texttt{NOT EXISTS}.
\end{block}
\end{frame}

\begin{frame}{Relational Algebra Cheat Sheet}
\small
\begin{center}
\renewcommand{\arraystretch}{1.25}
\begin{tabular}{lll}
\toprule
\textbf{Operator} & \textbf{Notation} & \textbf{SQL counterpart} \\
\midrule
Rename    & $\rho(R, r)$              & \texttt{FROM R r} \\
Selection & $\sigma_{[c]}(R)$         & \texttt{WHERE c} \\
Projection& $\pi_{[a_1,\ldots]}(R)$   & \texttt{SELECT DISTINCT a1, \ldots} \\
Cross product & $R \times S$          & \texttt{FROM R, S} \\
Join      & $R \bowtie_{[c]} S$       & \texttt{FROM R, S WHERE c} \\
Union / Difference & $R \cup S$, $R - S$ & \texttt{UNION}, \texttt{EXCEPT} \\
\bottomrule
\end{tabular}
\end{center}
\medskip
\begin{itemize}
  \item Relational algebra works on \textbf{sets}: $\pi$ removes duplicates automatically --- hence \texttt{DISTINCT} (or a set operator) in the SQL.
  \item $R \bowtie_{[c]} S \equiv \sigma_{[c]}(R \times S)$.
  \item $\cup$ and $-$ need \textbf{union-compatible} inputs (same number and types of columns).
\end{itemize}
\end{frame}

\section{Preliminary Questions}

\begin{frame}{Preliminary: Basic Relational Algebra}
\small
\textit{Not discussed in class --- attempt before the tutorial, since later questions build on these.}
\begin{enumerate}
  \item[\textbf{1.}] \textbf{Basic Relational Algebra.}
  \begin{enumerate}
    \item[(a)] Find the different departments in the School of Computing.
    \item[(b)] Print the emails of the students who \emph{borrowed} but did \emph{not lend} a copy of a book on the day that they joined the university.
  \end{enumerate}
\end{enumerate}
\end{frame}

\begin{frame}[fragile]{1(a). Departments in the School of Computing}
\small
Rename, select the faculty, project the department.
\[
\proj{d.department}\Big(\sel{d.faculty = 'School of Computing'}\big(\ren{department}{d}\big)\Big)
\]
\begin{lstlisting}[language=SQL]
SELECT d.department
FROM department d
WHERE d.faculty = 'School of Computing';
\end{lstlisting}
\texttt{department} is the primary key, so no duplicates can arise --- no \texttt{DISTINCT} needed.
\end{frame}

\begin{frame}[fragile]{1(b). Borrowed but Did Not Lend on Joining Day}
\footnotesize
``Borrowed on day one'' \textbf{minus} ``lent on day one'' --- a set difference.
\[
\begin{aligned}
&\proj{s1.email}\Big(\sel{s1.email = l1.borrower $\wedge$ l1.borrowed = s1.year}\big(\ren{student}{s1} \times \ren{loan}{l1}\big)\Big)\\
-\ &\proj{s2.email}\Big(\sel{s2.email = l2.owner $\wedge$ l2.borrowed = s2.year}\big(\ren{student}{s2} \times \ren{loan}{l2}\big)\Big)
\end{aligned}
\]
\begin{lstlisting}[language=SQL]
SELECT s1.email
FROM loan l1, student s1
WHERE s1.email = l1.borrower AND l1.borrowed = s1.year
EXCEPT
SELECT s2.email
FROM loan l2, student s2
WHERE s2.email = l2.owner AND l2.borrowed = s2.year;
\end{lstlisting}
\texttt{EXCEPT} is a set operator: it already removes duplicates.
\end{frame}

\section{Discussion Questions}

\begin{frame}{Discussion: Integrity Check \& Universal Quantification}
\small
\textit{To be discussed in class --- attempt before the tutorial; you may be called on to present.}
\begin{enumerate}
  \item[\textbf{2.}] \textbf{Relational Algebra.} Let us check the integrity of the data. Find the different emails of the students who borrowed \emph{or} lent a copy of a book \emph{before} they joined the university. There should not be any.
  \begin{enumerate}
    \item[(a)] Solve the problem \textbf{with} a set operation.
    \item[(b)] Solve the problem \textbf{without} a set operation.
  \end{enumerate}
  \item[\textbf{3.}] \textbf{Universal Quantification.} Print the emails and the names of the different students who borrowed \emph{all} the books authored by Adam Smith.
\end{enumerate}
\end{frame}

\begin{frame}[fragile]{2(a). Integrity Check --- With a Set Operation}
\footnotesize
``Borrowed too early'' \textbf{or} ``lent too early'' --- a union.
\[
\begin{aligned}
&\proj{s1.email}\Big(\sel{s1.email = l1.borrower $\wedge$ l1.borrowed $<$ s1.year}\big(\ren{student}{s1} \times \ren{loan}{l1}\big)\Big)\\
\cup\ &\proj{s2.email}\Big(\sel{s2.email = l2.owner $\wedge$ l2.borrowed $<$ s2.year}\big(\ren{student}{s2} \times \ren{loan}{l2}\big)\Big)
\end{aligned}
\]
\begin{lstlisting}[language=SQL]
SELECT s1.email
FROM loan l1, student s1
WHERE s1.email = l1.borrower AND l1.borrowed < s1.year
UNION
SELECT s2.email
FROM loan l2, student s2
WHERE s2.email = l2.owner AND l2.borrowed < s2.year;
\end{lstlisting}
Compare with 1(b): same shape, but $-$ becomes $\cup$ and $=$ becomes $<$.
\end{frame}

\begin{frame}[fragile]{2(b). Integrity Check --- Without a Set Operation}
\small
Fold the ``or'' into a \textbf{single selection} with a disjunction.
\[
\begin{aligned}
\proj{s.email}\Big(&\sel{(s.email = l.borrower $\vee$ s.email = l.owner) $\wedge$ l.borrowed $<$ s.year}\\
&\quad\big(\ren{student}{s} \times \ren{loan}{l}\big)\Big)
\end{aligned}
\]
\begin{lstlisting}[language=SQL]
SELECT DISTINCT s.email
FROM student s, loan l
WHERE (s.email = l.borrower OR s.email = l.owner)
  AND l.borrowed < s.year;
\end{lstlisting}
\begin{itemize}
  \item \textbf{\texttt{DISTINCT} is now required}: without a set operator, SQL keeps duplicates, but $\pi$ does not.
  \item \textbf{Mind the parentheses}: \texttt{AND} binds tighter than \texttt{OR}.
\end{itemize}
\end{frame}

\begin{frame}{3. Borrowed \emph{All} of Adam Smith's Books}
\small
In Tutorial 4 we wrote this with a \textbf{double \texttt{NOT EXISTS}}. But \texttt{NOT EXISTS} has no direct relational algebra counterpart --- we need to rephrase it with \textbf{set difference}.
\medskip

\textbf{Break it down:}
\begin{enumerate}
  \item Find the students for whom there is \textbf{at least one} Adam Smith book they have \emph{not} borrowed.
  \item Answer $=$ \textbf{all} students (including those who borrowed nothing) \textbf{minus} the students from Step 1.
\end{enumerate}
\medskip
\textbf{Step 1, split further:}
\begin{enumerate}
  \item[A.] Every \emph{possible} (student, Adam Smith book) pair --- a cross product.
  \item[B.] Every \emph{actual} (student, Adam Smith book) borrowing --- a join with \texttt{loan}.
  \item[C.] A $-$ B are the \textbf{missing} pairs; project to keep only email and name.
\end{enumerate}
\end{frame}

\begin{frame}{3. The Relational Algebra}
\scriptsize
We use $\coloneqq$ to name sub-expressions we refer to later.
\[
\begin{aligned}
Q_{1A} &\coloneqq \proj{s1.email, s1.name}\big(\ren{student}{s1}\big)\\
       &\qquad \times \proj{b1.ISBN13}\Big(\sel{b1.authors = 'Adam Smith'}\big(\ren{book}{b1}\big)\Big)\\[4pt]
Q_{1B} &\coloneqq \proj{s2.email, s2.name, b2.ISBN13}\Big(\ren{loan}{l2}\\
       &\qquad \njoin{l2.book = b2.ISBN13 $\wedge$ b2.authors = 'Adam Smith'} \ren{book}{b2}\\
       &\qquad \njoin{l2.borrower = s2.email} \ren{student}{s2}\Big)\\[4pt]
Q_{1} &\coloneqq \proj{s2.email, s2.name}\big(Q_{1A} - Q_{1B}\big)\\[4pt]
Q_{2} &\coloneqq \proj{s3.email, s3.name}\big(\ren{student}{s3}\big) - Q_1
\end{aligned}
\]
\medskip
\begin{itemize}
  \item $Q_{1A}$ and $Q_{1B}$ are \textbf{union-compatible}: both are (email, name, ISBN13).
  \item $Q_{1A} - Q_{1B}$: the (student, book) pairs that \emph{should} exist but \emph{don't}.
\end{itemize}
\end{frame}

\begin{frame}[fragile]{3. Translating Back to SQL}
\footnotesize
Each $-$ becomes an \texttt{EXCEPT}; $Q_1$ becomes a subquery in \texttt{FROM}.
\begin{lstlisting}[language=SQL,basicstyle=\ttfamily\scriptsize]
SELECT s3.email, s3.name
FROM student s3
EXCEPT
SELECT tmp.email, tmp.name
FROM (
  SELECT s1.email, s1.name, b1.ISBN13       -- Q1A
  FROM student s1, book b1
  WHERE b1.authors = 'Adam Smith'
  EXCEPT
  SELECT s2.email, s2.name, b2.ISBN13       -- Q1B
  FROM student s2, book b2, loan l2
  WHERE b2.authors = 'Adam Smith'
    AND s2.email = l2.borrower
    AND b2.ISBN13 = l2.book
) AS tmp;
\end{lstlisting}
Same answer as Tutorial 4's double-\texttt{NOT EXISTS} query: ``no Adam Smith book is missing from $s$'s loans.''
\end{frame}

\begin{frame}{3. Why This Is the Same as \texttt{NOT EXISTS}}
\small
\begin{center}
\renewcommand{\arraystretch}{1.3}
\begin{tabular}{ll}
\toprule
\textbf{Nested SQL (Tut 4)} & \textbf{Relational algebra (Tut 5)} \\
\midrule
inner \texttt{NOT EXISTS} (no loan of $b$ by $s$) & $Q_{1A} - Q_{1B}$ \\
outer \texttt{NOT EXISTS} (no such $b$ for $s$) & $\pi(\text{student}) - Q_1$ \\
\bottomrule
\end{tabular}
\end{center}
\medskip
\begin{itemize}
  \item $\forall b\,.\,\text{borrowed}(s,b) \;\equiv\; \neg\exists b\,.\,\neg\text{borrowed}(s,b)$
  \item Each \textbf{negation} becomes a \textbf{set difference} against the set of all candidates.
  \item Students who borrowed \emph{nothing} still appear in $Q_{1A}$ (via the cross product), so they correctly land in $Q_1$ and are excluded --- unless Adam Smith has no books, in which case \textbf{every} student qualifies (vacuous truth).
\end{itemize}
\end{frame}

\section{Guidelines}

\begin{frame}{Guidelines \& Common Pitfalls}
\small
\begin{itemize}\itemsep4pt
  \item \textbf{Sets vs.\ bags.} $\pi$ eliminates duplicates; a plain SQL \texttt{SELECT} does not. Add \texttt{DISTINCT} unless a set operator (or a key) already guarantees uniqueness (1(a), 2(b)).
  \item \textbf{``or'' $\to \cup$ or $\vee$; ``but not'' $\to -$.} Either a set operator between two queries, or a combined condition inside one $\sigma$ (Q2).
  \item \textbf{Parenthesise disjunctions.} \texttt{(A OR B) AND C} $\neq$ \texttt{A OR B AND C} (2(b)).
  \item \textbf{``For all'' has no single operator.} Build all candidates ($\times$), remove what actually happened ($-$), then remove the failures from everyone ($-$ again) (Q3).
  \item \textbf{Union-compatibility.} Both sides of $\cup$ / $-$ must have the same columns --- keep \texttt{ISBN13} in $Q_{1A}$ and $Q_{1B}$ \emph{until after} the difference.
\end{itemize}
\end{frame}

\begin{frame}
\begin{center}
Questions?\\
Drop a mail at: pratik.karmakar@u.nus.edu
\end{center}
\end{frame}

\end{document}
